% ============================================================================
% CONCLUSIONES
% ============================================================================

\section{Conclusiones}

\subsection{Hallazgos Principales}

\subsubsection{Construcción de la Tabla}

\begin{itemize}
    \item Base de 360,441 registros con 28,510 muertes (período 2000-2012)
    \item Graduación WH validada: Chi-cuadrado (p=0.9959), KS (p=0.9912), ajuste excelente
    \item Extrapolaciones:
    \begin{itemize}
        \item Cabeza: Oppermann (0-9 años) y Heligman-Pollard (10-19 años)
        \item Factor de empalme: $C = q_{10}^{HP} / q_{10}^{Oppermann}$
        \item Cola: Heligman-Pollard (85-109 años)
    \end{itemize}
    \item Mortalidad inferior a MI2014 en toda la tabla
\end{itemize}

\subsubsection{Valores de Conmutación}

\begin{itemize}
    \item Funciones calculadas correctamente ($D_x$, $N_x$, $S_x$, $C_x$, $M_x$)
    \item Relaciones recursivas verificadas, condiciones finales en $\omega=109$ satisfechas
    \item Valores sistemáticamente superiores vs. MI2014 (ej: $N_{25}$ +35\%)
\end{itemize}

\subsubsection{Valores Actuales}

\begin{itemize}
    \item Rentas diferidas: diferenciales hasta +115\% promedio
    \item Dotal puro: +58\% promedio
    \item Rentas temporales: +18\% promedio
    \item Tabla obtenida genera valores sistemáticamente mayores (14 de 15 casos)
    \item MI2014 subvalúa productos de supervivencia para esta población
\end{itemize}

\subsection{Implicaciones Actuariales}

\begin{itemize}
    \item \textbf{Primas:} MI2014 subestima costo de productos de supervivencia (+18\% a +115\% según producto)
    \item \textbf{Reservas:} Necesidad de ajuste patrimonial para rentas vitalicias
    \item \textbf{Riesgo de longevidad:} Mayor supervivencia requiere provisiones adicionales
    \item \textbf{Actualización:} Revisión periódica (3-5 años) de tablas recomendada
\end{itemize}

\subsubsection{Para la Política de Productos}

Los resultados orientan decisiones estratégicas sobre cartera de productos:

\begin{itemize}
    \item \textbf{Sensibilidad diferencial:} Las rentas diferidas (diferencial +115\%) y dotales puros (+58\%) son extremadamente sensibles a diferencias de mortalidad. Requieren precaución especial en tarifación y gestión de cartera.
    
    \item \textbf{Mitigación de riesgo:} Estrategias recomendadas incluyen: (1) reaseguro selectivo en productos de mayor sensibilidad, (2) límites de concentración en rentas vitalicias, (3) cláusulas de revisión de primas basadas en experiencia observada.
    
    \item \textbf{Pricing diferenciado:} Implementar tarifación segmentada por tipo de invalidez y características demográficas, reconociendo la heterogeneidad de mortalidad dentro de la población de inválidos.
    
    \item \textbf{Rentabilidad por línea:} Los productos temporales (diferencial +18\%) presentan menor sensibilidad y riesgo, siendo más predecibles en rentabilidad. Las rentas vitalicias requieren márgenes de seguridad mayores por su exposición prolongada al riesgo de longevidad.
\end{itemize}

\subsection{Limitaciones}

\begin{itemize}
    \item Exposición limitada en edades extremas (0-19 y 85+) requiere extrapolaciones
    \item Datos del período 2000-2012 podrían no reflejar mejoras recientes
    \item Especificidad poblacional: resultados aplicables a población analizada
    \item Método WH asume suavidad, puede enmascarar discontinuidades reales
\end{itemize}
\subsection*{Anexos}

\textbf{Material complementario disponible en el repositorio del proyecto:}

\begin{itemize}
    \item \texttt{trabajo\_final.r} - Script principal de análisis
    \item \texttt{tabla\_mortalidad\_obtenida.csv} - Tabla completa construida
    \item \texttt{tasas\_crudas.csv} - Tasas brutas calculadas
    \item \texttt{tabla\_referencia\_MI2014.csv} - Tabla de referencia
    \item \texttt{datos\_graduacion.csv} - Datos del proceso de graduación
    \item Carpeta \texttt{imagenes/} - Gráficos y visualizaciones generadas
\end{itemize}

\textbf{Paquetes de R utilizados:}
\begin{itemize}
    \item Wickham H, et al. (2019). \textit{tidyverse: Easily Install and Load the 'Tidyverse'}. R package version 1.3.0.
    \item Grolemund G, Wickham H (2011). \textit{lubridate: Make Dealing with Dates a Little Easier}. R package.
    \item \textit{MortalityTables}: Tablas de mortalidad y graduación. R package.
    \item \textit{MortalityLaws}: Modelos de mortalidad y extrapolación. R package.
    \item \textit{randtests}: Tests estadísticos no paramétricos. R package.
    \item \textit{BSDA}: Basic Statistics and Data Analysis. R package.
    \item \textit{tseries}: Time Series Analysis and Computational Finance. R package.
\end{itemize}

\textbf{Nota sobre uso de IA:} Parte del código de este proyecto fue asistido mediante herramientas de inteligencia artificial para optimización y depuración.

% ----------------------------------------------------------------------------
\vspace{1cm}

\begin{center}
\begin{tcolorbox}[colback=celestefondo, colframe=celesteprincipal, boxrule=1.5pt, arc=4pt, width=0.9\textwidth]
\centering
\textcolor{celesteoscuro}{\textbf{Fin del Informe}}

\vspace{0.3cm}

\textcolor{grisoscuro}{\textit{Proyecto Final - EYP2605 Matemáticas Actuariales}}

\textcolor{grisoscuro}{\textit{Pontificia Universidad Católica de Chile}}

\vspace{0.2cm}

\textcolor{celesteprincipal}{Diciembre 2020}
\end{tcolorbox}
\end{center}
